\documentclass[12pt, a4paper]{article}

\usepackage[margin=1in]{geometry}
\usepackage{graphicx}
\usepackage{amsmath}
\usepackage{booktabs}
\usepackage{enumitem}
\usepackage{hyperref}
\usepackage{setspace}
\usepackage{parskip}
\usepackage{longtable}
\usepackage{float}

\onehalfspacing

\title{\textbf{Structural Thinking Framework for Solving Data Science Problems} \\ \large Applied to the EV Charging Dataset}
\author{[Your Name Here] \\ [Your Registration Number] \\ VIT University}
\date{\today}

\begin{document}

\maketitle

\tableofcontents
\newpage

% ============================================================
\section{Step 1: Clarify the Problem}
% ============================================================

Whenever we start working on a data science project, the first thing we need to do is actually understand what we're trying to solve. It sounds really obvious, but honestly a lot of times people just jump straight into coding and plotting graphs without even knowing what question they are actually answering. I've done this too, and it usually ends up wasting a lot of time.

So for our EV Charging dataset, we need to sit down and figure out a few things first.

\subsection{What is the real question being asked?}

The dataset we have contains information about electric vehicle charging sessions -- stuff like when the charging happened, how long it took, how much energy got consumed, what kind of charger was used, and so on. But just having data doesn't mean much unless we know what we want to find out from it.

For us, the main question could be something like: \textit{``Can we predict the energy consumption or charging cost of an EV charging session based on the available features?''} Or maybe we want to understand the \textit{``What factors influence charging behaviour the most?''} -- both are valid questions, but we need to pick one and be clear about it.

\subsection{What decision will this analysis help make?}

This is about thinking from the perspective of whoever is going to use our results. Like if you're running a charging station network, maybe you want to know which stations are getting overloaded and when, so that you can plan better. Or if you're an EV manufacturer, you might want to understand how people charge their vehicles so you can design better batteries.

In our case, lets say the analysis is meant to help a charging network operator decide things like -- where to put new stations, what charger types to invest in, how to set pricing, etc.

\subsection{What is the success metric?}

We need to define upfront how we'll know if our model or analysis is actually working well. For a prediction problem, this might be something like RMSE (Root Mean Squared Error) or $R^2$ score. For a classification task, we'd look at accuracy or F1-score.

For our EV dataset, if we go with predicting charging cost, then our success metric could be something like:
\begin{itemize}
    \item RMSE $\leq$ some threshold (say $\delta$ USD)
    \item $R^2 > 0.7$ or so
\end{itemize}

We don't have to be super precise about these numbers right now, but having a rough target in mind really helps keep things on track.

\subsection{What outcome does the stakeholder expect?}

The stakeholder (in our case, lets just assume its our professor or a charging network company) probably expects something like -- a clear analysis showing which factors affect charging patterns, and maybe a working model that can predict costs or energy usage with reasonable accuracy. They'd also want some visualisations and a summary of findings that actually make sense to someone who isn't a data scientist.

\textbf{Goal:} We convert our vague problem into a clear, measurable statement. Something like: \textit{``Predict the Charging Cost (USD) for a given EV charging session with an $R^2$ score above 0.7, using features such as battery capacity, charging duration, charger type, etc.''}

\newpage

% ============================================================
\section{Step 2: Decompose the Problem}
% ============================================================

Now that we have a clear problem statement, the next step is to break it down into smaller pieces. Big problems are always scary, but when you break them into manageable chunks, they become much easier to handle.

For our EV charging dataset, the decomposition looks something like this:

\subsection{Sub-problem 1: Define the target variable}

We need to clearly define what we're trying to predict or analyze. In our case, the target variable is \texttt{Charging Cost (USD)}. But we also need to check -- is this column clean? Does it have missing values? Are there any weird outliers where the cost is 0 or negative?

\subsection{Sub-problem 2: Identify influencing factors}

Which columns in the dataset could possibly affect the charging cost? Some obvious candidates are:
\begin{itemize}
    \item Energy Consumed (kWh) -- more energy used = higher cost, probably
    \item Charging Duration (hours) -- longer sessions might cost more
    \item Charger Type -- DC Fast Chargers are usually pricier
    \item Charging Rate (kW)
    \item Battery Capacity (kWh) -- bigger battery might mean bigger session
    \item Location -- prices might vary across cities
\end{itemize}

But there could also be less obvious ones like Temperature or Time of Day that affect things indirectly. We wont know for sure until we explore the data.

\subsection{Sub-problem 3: Choose the prediction approach}

Do we want a simple regression model, or something more complex like a random forest or gradient boosting? We'll decide this later, but for now lets just note that this is a regression problem since our target is a continuous variable (cost in USD).

\subsection{Sub-problem 4: Decide on performance metrics}

As discussed in Step 1, we'll use RMSE, MAE, and $R^2$ as our main metrics. But we should also think about things like -- does the model perform equally well for cheap sessions vs expensive ones? Are there specific user types where the model fails badly?

This decomposition reduces the overall complexity and makes sure we don't miss anything important.

\newpage

% ============================================================
\section{Step 3: Identify What Data Is Needed}
% ============================================================

For each sub-problem identified in Step 2, we need to figure out what data we actually need and whether our dataset has it.

\subsection{Data source assessment}

Luckily, we already have a dataset -- the EV Charging Dataset with around 1300+ records and 20 columns. Lets list out what we have:

\begin{table}[H]
\centering
\caption{Dataset columns and their types}
\begin{tabular}{p{5.5cm} p{3cm} p{4cm}}
\toprule
\textbf{Column Name} & \textbf{Type} & \textbf{Role} \\
\midrule
User ID & Categorical & Identifier \\
Vehicle Model & Categorical & Feature \\
Battery Capacity (kWh) & Numerical & Feature \\
Charging Station ID & Categorical & Identifier \\
Charging Station Location & Categorical & Feature \\
Charging Start Time & DateTime & Feature \\
Charging End Time & DateTime & Feature \\
Energy Consumed (kWh) & Numerical & Feature / Target \\
Charging Duration (hours) & Numerical & Feature \\
Charging Rate (kW) & Numerical & Feature \\
Charging Cost (USD) & Numerical & Target \\
Time of Day & Categorical & Feature \\
Day of Week & Categorical & Feature \\
State of Charge - Start (\%) & Numerical & Feature \\
State of Charge - End (\%) & Numerical & Feature \\
Distance Driven (km) & Numerical & Feature \\
Temperature (\textdegree C) & Numerical & Feature \\
Vehicle Age (years) & Numerical & Feature \\
Charger Type & Categorical & Feature \\
User Type & Categorical & Feature \\
\bottomrule
\end{tabular}
\end{table}

\subsection{Do we have historical data?}

Yes, the dataset covers charging sessions starting from January 2024 and has timestamps. So we do have temporal information which could be usefull for spotting trends or seasonal patterns.

\subsection{Data quality expectations}

Before we even start analysis, we should expect that there will be some issues:
\begin{itemize}
    \item \textbf{Missing values:} Some columns like Energy Consumed, Charging Rate, and Distance Driven might have NaN entries. We noticed that some rows have empty fields in the CSV.
    \item \textbf{Inconsistent data:} Battery capacity values seem to vary a lot -- some are standard values like 50, 62, 75, 85, 100 kWh but others look like non-standard decimal values (e.g. 108.46, 48.79) which could be noise or actual variant models.
    \item \textbf{Potential outliers:} State of Charge values going above 100\% or below 0\% (we spotted values like 132.9\%, 152.4\%, -9.3\textdegree C temperature which could be erroneous or just extreme weather).
    \item \textbf{Data types:} Start and End times are stored as strings, so we'll need to convert them to proper datetime format.
\end{itemize}

\textbf{Goal:} A complete list of required data attributes -- which we have, thankfully. We just need to clean and preprocess everything properly before we can use it.

\newpage

% ============================================================
\section{Step 4: Form Hypotheses}
% ============================================================

This is actually one of my favourite steps because you get to think about the problem like a detective before looking at the evidence. The idea is to come up with some educated guesses about what patterns might exist in the data. Then later during EDA, we can check if these guesses actually hold up.

Here are the hypotheses I came up with for our EV charging dataset:

\begin{enumerate}
    \item \textbf{H1:} ``Higher energy consumption leads to higher charging costs.'' \\
    This seems pretty straightforward -- if you use more electricity, you pay more. But the relationship might not be perfectly linear depending on the pricing structure.

    \item \textbf{H2:} ``DC Fast Chargers cost more per session compared to Level 1 and Level 2 chargers.'' \\
    Fast charging is usually more expensive because of the infrastructure costs and the convenience factor.

    \item \textbf{H3:} ``Charging sessions during peak hours (Afternoon/Evening) are more expensive than off-peak (Night/Morning).'' \\
    This is based on the idea that electricity pricing often follows time-of-use tariffs.

    \item \textbf{H4:} ``Longer charging duration does not necessarily mean higher cost.'' \\
    Because it depends on the charging rate -- a slow charger running for 4 hours might deliver less energy than a fast charger running for 1 hour.

    \item \textbf{H5:} ``Commuters have different charging patterns compared to Casual Drivers and Long-Distance Travelers.'' \\
    Different user types probably have different needs and hence different charging behaviours.

    \item \textbf{H6:} ``Temperature has some effect on charging efficiency or energy consumed.'' \\
    This is based on the known fact that EV battery performance is affected by ambient temperature -- batteries are less efficient in extreme cold or heat.

    \item \textbf{H7:} ``Charging cost varies significantly across different cities (locations).'' \\
    Electricity prices differ from city to city, so Houston might be cheaper than San Francisco, for example.
\end{enumerate}

These hypotheses give us a clear direction for our exploratory data analysis. Instead of just making random plots, we now have specific things to look for. Some of these might turn out to be wrong, and thats completely fine -- thats how data science works.

\newpage

% ============================================================
\section{Step 5: Select the Right Analytical Approach}
% ============================================================

Now we need to decide what kind of analytical method to use. The choice depends entirely on what type of problem we're solving. Heres a quick summary of the common problem types and their corresponding methods:

\begin{table}[H]
\centering
\caption{Problem type vs. Method mapping}
\begin{tabular}{p{5cm} p{6cm}}
\toprule
\textbf{Problem Type} & \textbf{Method} \\
\midrule
Predict a number & Regression \\
Predict a category & Classification \\
Segment groups & Clustering \\
Detect unusual events & Anomaly Detection \\
Forecast future values & Time Series Forecasting \\
\bottomrule
\end{tabular}
\end{table}

\subsection{Our approach}

Since our target variable (Charging Cost in USD) is continuous/numerical, this is clearly a \textbf{regression problem}. We want to predict a number, not a category.

However, we can also do some supplementary analysis:
\begin{itemize}
    \item \textbf{Clustering:} We could cluster users based on their charging patterns to identify distinct user segments (even though we already have User Type, there might be more nuanced groupings).
    \item \textbf{Anomaly Detection:} We could flag unusual charging sessions -- like ones where the cost is way too high for the energy consumed, or where the SoC goes above 100\%.
\end{itemize}

For the main regression task, we plan to try multiple algorithms:
\begin{enumerate}
    \item Linear Regression (baseline)
    \item Decision Tree Regressor
    \item Random Forest Regressor
    \item Gradient Boosting (XGBoost or similar)
\end{enumerate}

We'll compare their performance using the metrics defined earlier ($R^2$, RMSE, MAE) and pick the best one. This gives us a methodological blueprint -- we know what we're doing and why before we write a single line of model training code.

\newpage

% ============================================================
\section{Step 6: Prepare and Explore the Data}
% ============================================================

This is where the actual hands-on work begins. Data preparation is probably the most time-consuming part of any data science project -- they say data scientists spend 80\% of their time just cleaning data, and honestly I believe it.

\subsection{Data Preparation}

\subsubsection{Handling Missing Values}

The first thing we do is check for missing values across all columns. Based on our initial look at the CSV, several columns have missing entries:

Let $N$ be the total number of rows and $n_{\text{miss},j}$ be the count of missing values in column $j$. We calculate the missing percentage as:
$$\text{Missing \%}_j = \frac{n_{\text{miss},j}}{N} \times 100$$

For columns with low missing percentage (say $< 5\%$), we can either drop those rows or impute them using mean/median. For example:
\begin{itemize}
    \item For numerical columns like \texttt{Energy Consumed} or \texttt{Charging Rate} -- use median imputation (since mean can be skewed by outliers)
    \item For categorical columns -- use mode imputation
\end{itemize}

If a column has too many missing values (say $> 30\%$), we might consider dropping it entirely. But from what we observed, none of our columns seem to be that bad.

\subsubsection{Removing Outliers}

Outliers can seriously mess up our models. We'll use the IQR (Interquartile Range) method to detect them:

$$\text{IQR} = Q_3 - Q_1$$
$$\text{Lower Bound} = Q_1 - 1.5 \times \text{IQR}$$
$$\text{Upper Bound} = Q_3 + 1.5 \times \text{IQR}$$

Any values falling outside $[\text{Lower Bound}, \text{Upper Bound}]$ are flagged as potential outliers. For our dataset, we need to check columns like:
\begin{itemize}
    \item \texttt{Battery Capacity} -- values like 3.97, 1.53 kWh seem unrealistically low
    \item \texttt{State of Charge (Start/End)} -- values above 100\% or below 0\% are problematic
    \item \texttt{Temperature} -- values like $-9.8$\textdegree C might be valid for some cities but worth checking
    \item \texttt{Vehicle Age} -- values above 7-8 years for EVs could be suspicious
\end{itemize}

\subsubsection{Feature Engineering}

This is where we create new columns from existing data that might be more useful for our model:

\begin{itemize}
    \item \texttt{Charge Duration} -- we can validate/recalculate this from the difference between end time and start time: $\Delta t = t_{\text{end}} - t_{\text{start}}$
    \item \texttt{SoC Change} -- calculate the change in state of charge: $\Delta \text{SoC} = \text{SoC}_{\text{end}} - \text{SoC}_{\text{start}}$
    \item \texttt{Hour of Day} -- extract from the start time for more granular time analysis
    \item \texttt{Is Weekend} -- binary flag: 1 if Saturday/Sunday, 0 otherwise
    \item \texttt{Cost per kWh} -- derived metric: $\frac{\text{Charging Cost}}{\text{Energy Consumed}}$
\end{itemize}

\subsubsection{Encoding and Scaling}

Categorical variables like Vehicle Model, Charger Type, User Type, Time of Day, etc. need to be converted to numbers before we can feed them to most ML models.

\begin{itemize}
    \item \textbf{Label Encoding:} For ordinal categories (if any)
    \item \textbf{One-Hot Encoding:} For nominal categories like Vehicle Model, Location, Charger Type
\end{itemize}

For numerical features, we apply standardisation:
$$x_{\text{scaled}} = \frac{x - \mu}{\sigma}$$

where $\mu$ is the mean and $\sigma$ is the standard deviation. This ensures all features are on a similar scale, which is important for algorithms like linear regression that are sensitive to feature magnitude.

\subsubsection{Train-Test Split}

We split the data into training and testing sets. A standard 80-20 split works fine:

$$D = D_{\text{train}} \cup D_{\text{test}}, \quad |D_{\text{train}}| = 0.8 \times |D|, \quad |D_{\text{test}}| = 0.2 \times |D|$$

We use \texttt{random\_state = 42} (or some fixed seed) so that our results are reproducible. Its also worth considering stratified splitting if our target variable has a very uneven distribution, but for a continuous variable we usually just do random splitting.

\subsection{Exploratory Data Analysis (EDA)}

This is the fun part -- actually looking at the data through plots and statistics to understand its story.

\subsubsection{Identify Patterns}

Some of the things we look for:
\begin{itemize}
    \item Distribution of charging costs -- is it normal? Skewed? Has multiple peaks?
    \item How charging cost relates to energy consumed (scatter plot)
    \item Average cost by charger type (bar chart)
    \item Charging sessions by time of day (histogram)
    \item Geographic distribution -- which city has the most sessions?
\end{itemize}

\subsubsection{Validate Hypotheses}

We go back to our hypotheses from Step 4 and check them one by one:
\begin{itemize}
    \item For H1, we'd plot Energy Consumed vs Charging Cost and compute the Pearson correlation coefficient $r$. If $r$ is close to 1, our hypothesis is supported.
    \item For H2, we'd compute the mean cost grouped by Charger Type and do an ANOVA test or just compare the group means visually.
    \item For H3, we group by Time of Day and compare mean costs.
    \item And so on for each hypothesis...
\end{itemize}

\subsubsection{Understand Correlations}

We compute a correlation matrix for all numerical columns:
$$\rho_{ij} = \frac{\text{Cov}(X_i, X_j)}{\sigma_{X_i} \cdot \sigma_{X_j}}$$

This helps us identify which features are strongly correlated with our target (good for modeling) and which features are correlated with each other (potential multicollinearity issues).

\subsubsection{Identify Data Limitations}

During EDA, we'll probably discover some limitations:
\begin{itemize}
    \item The dataset might be synthetic or simulated (some values seem too perfectly distributed)
    \item All sessions seem to be from 2024, so temporal generalization might be limited
    \item We only have 5-6 cities, so geographic coverage is limited
    \item Some combinations of features might not make physical sense (e.g., very high energy consumed from a very low battery capacity)
\end{itemize}

Being honest about these limitations is important because it affects how much we can trust our final results.

\newpage

% ============================================================
\section{Step 7: Model Building}
% ============================================================

Alright, so this is where we actually build the machine learning models. By now we've cleaned the data, done our EDA, and have a pretty good understanding of what we're working with. Time to let the algorithms do their thing.

\subsection{Baseline Model}

We always start with a simple baseline. For regression, the simplest baseline is just predicting the mean value for every sample:

$$\hat{y}_{\text{baseline}} = \bar{y} = \frac{1}{n} \sum_{i=1}^{n} y_i$$

This gives us a ``floor'' to beat. If our actual model can't even beat the mean prediction, then something is seriously wrong. The RMSE of the baseline model would be approximately equal to the standard deviation of the target variable.

Next, we try a simple \textbf{Linear Regression}:

$$\hat{y} = \beta_0 + \beta_1 x_1 + \beta_2 x_2 + \ldots + \beta_p x_p$$

where $x_1, x_2, \ldots, x_p$ are our features and $\beta_0, \beta_1, \ldots, \beta_p$ are the coefficients learned from the training data.

\subsection{Advanced Models}

After the baseline, we move to more sophisticated models:

\begin{enumerate}
    \item \textbf{Decision Tree Regressor:} Simple, interpretable, but prone to overfitting. It splits the feature space into regions and predicts the mean of each region.
    
    \item \textbf{Random Forest Regressor:} An ensemble of decision trees. Each tree is trained on a random subset of the data and features, and the final prediction is the average of all trees:
    $$\hat{y}_{\text{RF}} = \frac{1}{B} \sum_{b=1}^{B} \hat{f}_b(x)$$
    where $B$ is the number of trees.
    
    \item \textbf{Gradient Boosting Regressor:} Similar to random forest but trees are built sequentially, each one trying to correct the errors of the previous one. Generally gives better performance but is more complex to tune.
\end{enumerate}

\subsection{Hyperparameter Tuning}

Each model has parameters that we need to set before training (hyperparameters). We use techniques like:

\begin{itemize}
    \item \textbf{Grid Search:} Try all combinations of a predefined parameter grid
    \item \textbf{Random Search:} Randomly sample combinations (faster for large grids)
    \item \textbf{Cross-Validation:} We use k-fold CV (usually $k=5$) to evaluate each parameter combination
\end{itemize}

For example, for Random Forest, we might tune:
\begin{itemize}
    \item \texttt{n\_estimators} $\in \{50, 100, 200\}$
    \item \texttt{max\_depth} $\in \{5, 10, 15, \texttt{None}\}$
    \item \texttt{min\_samples\_split} $\in \{2, 5, 10\}$
\end{itemize}

\subsection{Performance Comparison}

After training all models, we create a comparison table. The format would look something like:

\begin{table}[H]
\centering
\caption{Model Performance Comparison (notional)}
\begin{tabular}{lccc}
\toprule
\textbf{Model} & \textbf{$R^2$} & \textbf{RMSE} & \textbf{MAE} \\
\midrule
Mean Baseline & $\approx 0$ & $\sigma_y$ & -- \\
Linear Regression & $r_1$ & $e_1$ & $m_1$ \\
Decision Tree & $r_2$ & $e_2$ & $m_2$ \\
Random Forest & $r_3$ & $e_3$ & $m_3$ \\
Gradient Boosting & $r_4$ & $e_4$ & $m_4$ \\
\bottomrule
\end{tabular}
\end{table}

(Where $r_i$, $e_i$, $m_i$ denote the respective metric scores for each model.)

\textbf{Goal:} At the end of this step, we have the best-performing model that meets (or hopefully exceeds) our business objectives from Step 1.

\newpage

% ============================================================
\section{Step 8: Evaluate the Results}
% ============================================================

Building a model is one thing, but knowing if its actually good enough is another. This step is about rigorously evaluating our models and making sure the results are trustworthy.

\subsection{Using the Right Metrics}

Since this is a regression problem, we use:

\begin{itemize}
    \item \textbf{$R^2$ (Coefficient of Determination):} 
    $$R^2 = 1 - \frac{\sum_{i=1}^{n}(y_i - \hat{y}_i)^2}{\sum_{i=1}^{n}(y_i - \bar{y})^2}$$
    This tells us how much of the variance in the target is explained by our model. Ranges from 0 to 1 (higher is better), but can be negative for very bad models.
    
    \item \textbf{RMSE (Root Mean Squared Error):}
    $$\text{RMSE} = \sqrt{\frac{1}{n} \sum_{i=1}^{n}(y_i - \hat{y}_i)^2}$$
    This gives us the average magnitude of the prediction errors, in the same units as the target (USD in our case).
    
    \item \textbf{MAE (Mean Absolute Error):}
    $$\text{MAE} = \frac{1}{n} \sum_{i=1}^{n} |y_i - \hat{y}_i|$$
    Similar to RMSE but less sensitive to large errors because it doesn't square them.
\end{itemize}

If we had also done a classification task (like predicting User Type), we would use metrics like Accuracy, F1-score, and ROC-AUC. And for clustering, something like Silhouette Score.

\subsection{Is the model reliable?}

Reliability is about consistency. We check this through:
\begin{itemize}
    \item \textbf{Cross-validation scores:} If the model gives wildly different $R^2$ values across different folds, it might not be reliable.
    \item \textbf{Training vs Testing performance gap:} If training $R^2$ is 0.99 but testing $R^2$ is 0.40, the model is overfitting badly.
    \item \textbf{Residual analysis:} Plot residuals ($y - \hat{y}$) against predicted values. Ideally, residuals should be randomly scattered around zero with no obvious pattern.
\end{itemize}

\subsection{Does it generalize well?}

To test generalisation, we look at performance on the held-out test set. If the test performance is close to the training performance (within a reasonable margin), the model generalises well.

Mathematically, we want:
$$|R^2_{\text{train}} - R^2_{\text{test}}| < \epsilon$$

for some small $\epsilon$ (say 0.05-0.10). A large gap indicates overfitting.

\subsection{Does it meet the business success metric?}

Going back to Step 1 where we defined our success criteria -- if we wanted $R^2 > 0.7$ and RMSE below a certain threshold, we check whether our best model actually achieves that. If it doesn't, we might need to go back to Steps 6-7 and try different approaches -- maybe better feature engineering, more data cleaning, or a different algorithm altogether.

\newpage

% ============================================================
\section{Step 9: Interpret and Translate Insights}
% ============================================================

Having a model with good numbers is great, but its meaningless if we cant explain what it actually means in plain English. This is especially important because most people who'll use our results (managers, operators, stakeholders) don't care about $R^2$ scores -- they care about actionable information.

\subsection{Identify Key Features}

Most ML models can tell us which features were most important for making predictions. For tree-based models like Random Forest or Gradient Boosting, we get \textbf{feature importance scores}. The feature importance $I_j$ for feature $j$ represents how much that feature contributed to reducing the prediction error.

For our EV charging dataset, the top features influencing charging cost might turn out to be something like:
\begin{enumerate}
    \item Energy Consumed (kWh) -- $I_1$: probably the highest
    \item Charger Type -- $I_2$
    \item Charging Duration -- $I_3$
    \item Location -- $I_4$
    \item Charging Rate -- $I_5$
\end{enumerate}

(The exact rankings and importance values would come from the actual model.)

\subsection{Connect to Real-World Behaviour}

This is about translating numbers into stories:
\begin{itemize}
    \item ``Energy consumption is the single biggest driver of cost'' $\rightarrow$ this makes physical sense because you're paying for the electricity
    \item ``DC Fast Chargers have on average $\Delta$ USD higher cost per session'' $\rightarrow$ this aligns with industry practice where convenience charging costs more
    \item ``Evening sessions tend to be costlier'' $\rightarrow$ could be due to peak pricing or higher demand
    \item ``Newer vehicles tend to have slightly different charging patterns'' $\rightarrow$ newer EVs might have better battery management systems
\end{itemize}

\subsection{Provide Clear Explanations for Non-Technical Stakeholders}

Instead of saying \textit{``the model achieved an $R^2$ of 0.82 with 15 features''}, we could say:

\textit{``Our model can predict the charging cost of a session with about 82\% accuracy, based primarily on how much energy is consumed, what type of charger is used, and where the charging station is located. This means that for most sessions, our prediction will be within a few dollars of the actual cost.''}

This kind of translation is really important and honestly its a skill that a lot of data science students (including me) need to develop better. Numbers and formulas are great, but if you cant communicate your findings effectively, it doesnt matter how good your model is.

\newpage

% ============================================================
\section{Step 10: Recommend Actions}
% ============================================================

This step is about going from ``what did we find'' to ``what should we do about it.'' Its basically converting our insights into concrete business recommendations.

\subsection{What should the business do?}

Based on our analysis of the EV charging dataset, here are some potential recommendations:

\begin{enumerate}
    \item \textbf{Pricing Strategy:} If DC Fast Chargers are significantly more expensive, the operator could consider dynamic pricing -- maybe offer discounts during off-peak hours to encourage usage distribution throughout the day. If $\text{Cost}_{\text{peak}} >> \text{Cost}_{\text{off-peak}}$, there might be an opportunity for a pricing model based on demand.
    
    \item \textbf{Station Placement:} If certain locations consistently show higher demand or longer queue times, those areas need more stations. The analysis of charging sessions by location can inform expansion decisions.
    
    \item \textbf{Charger Type Investment:} If users are increasingly preferring DC Fast Chargers (even at higher costs), the company should invest more in fast charging infrastructure. But if Level 2 chargers have better utilisation rates, maybe a mixed approach works best.
    
    \item \textbf{User Segmentation:} Different user types (Commuter, Casual Driver, Long-Distance Traveler) have different needs. The company could create tailored subscription plans -- like a monthly plan for commuters who charge daily, and a pay-per-use model for casual drivers.
\end{enumerate}

\subsection{How will it impact KPIs?}

If implemented properly:
\begin{itemize}
    \item Dynamic pricing could increase revenue by encouraging off-peak usage
    \item Better station placement could reduce wait times and improve customer satisfaction
    \item Targeted marketing to different user segments could improve retention
\end{itemize}

We can estimate the expected impact using our model -- for instance, if reducing prices by $x\%$ during night hours increases session count by $y\%$, is the net revenue positive?

\subsection{What can be automated?}

Some things that could be automated using our model:
\begin{itemize}
    \item Real-time cost estimation for users (``Your session will cost approximately \$X'')
    \item Anomaly alerts for unusual charging sessions (potential fraud or equipment issues)
    \item Demand forecasting for load management
\end{itemize}

\subsection{Risks and Limitations}

We need to be honest about what could go wrong:
\begin{itemize}
    \item The model is trained on this specific dataset which may not represent all real-world scenarios
    \item External factors like government subsidies, electricity price changes, or new EV models could change the patterns significantly
    \item If the dataset is synthetic (which some values suggest), the recommendations might not directly apply to real operations
    \item Model predictions are estimates, not guarantees -- there will always be some error margin
\end{itemize}

\newpage

% ============================================================
\section{Step 11: Deploy and Monitor (For ML Projects)}
% ============================================================

If this was a real-world project where our model actually needs to run in production, we'd need to think about deployment and monitoring. Even for an academic project, its good to at least know what this step looks like.

\subsection{Deploy the Model}

Once we've selected our best model (say its the Gradient Boosting model with tuned hyperparameters), we need to deploy it so it can make predictions on new, incoming data.

The deployment pipeline would look something like:

\begin{enumerate}
    \item \textbf{Save the model:} Export the trained model to a file using something like \texttt{joblib} or \texttt{pickle} in Python.
    $$\texttt{joblib.dump(model, `ev\_cost\_predictor.pkl')}$$
    
    \item \textbf{Create an API:} Build a REST API (using Flask or FastAPI) that accepts input features as a JSON request, runs them through the model, and returns the predicted cost.
    
    \item \textbf{Host it:} Deploy the API on a cloud service (AWS, GCP, Azure) or even a simple server for testing.
    
    \item \textbf{Integrate:} Connect the API with the charging station's existing software system so that predictions can be made in real-time.
\end{enumerate}

\subsection{Set Up Pipelines}

A data pipeline ensures that:
\begin{itemize}
    \item New charging session data is automatically collected and stored
    \item Data preprocessing (cleaning, encoding, scaling) happens automatically before prediction
    \item Predictions are logged for later analysis
\end{itemize}

This is usually done using tools like Apache Airflow, Kubeflow, or simple cron jobs for smaller projects.

\subsection{Monitor Drift, Decay, and Errors}

Over time, models can get worse. This happens because of:

\begin{itemize}
    \item \textbf{Data Drift:} The distribution of incoming data changes. For example, if a new EV model with different battery characteristics becomes popular, the model might not handle it well.
    $$P(\mathbf{X}_{\text{new}}) \neq P(\mathbf{X}_{\text{train}})$$
    
    \item \textbf{Concept Drift:} The relationship between features and target changes. Maybe electricity prices change, or the company introduces new pricing tiers.
    $$P(Y|\mathbf{X})_{\text{new}} \neq P(Y|\mathbf{X})_{\text{train}}$$
    
    \item \textbf{Prediction Errors:} Tracking actual vs predicted values over time. If the error starts increasing consistently, its a sign the model needs retraining.
\end{itemize}

We set up monitoring dashboards (using tools like Grafana or even simple logging) to track these metrics over time.

\subsection{Retrain Periodically}

Based on the monitoring results, we retrain the model periodically -- maybe every month or every quarter, depending on how fast things change. The retraining process would use the latest data and follow the same pipeline as the original training (Steps 6-8).

A common strategy is to set threshold alerts:
$$\text{If } R^2_{\text{current}} < R^2_{\text{threshold}} \Rightarrow \text{Trigger retraining}$$

\textbf{Goal:} Stable, reliable, ongoing performance of our prediction model in a real production environment.

\newpage

% ============================================================
\section{Conclusion}
% ============================================================

To summarize, the 11-step structural thinking framework provides a systematic way to approach any data science problem. Instead of jumping straight into coding (which is what most of us naturally want to do), it forces us to think carefully about the problem, plan our approach, and execute methodically.

For our EV Charging dataset, following this framework helped us:

\begin{enumerate}
    \item Clearly define what we want to predict (Charging Cost)
    \item Break the problem into manageable sub-tasks
    \item Identify what data we need and what we already have
    \item Form testable hypotheses about charging patterns
    \item Choose regression as our main analytical approach
    \item Clean and explore the data thoroughly
    \item Build and compare multiple models
    \item Evaluate them using appropriate metrics
    \item Translate findings into understandable insights
    \item Make actionable business recommendations
    \item Plan for deployment and monitoring
\end{enumerate}

The key takeaway is that data science is not just about building fancy models -- its about solving real problems in a structured, thoughtful way. Every step in this framework serves a purpose, and skipping any of them can lead to incomplete or unreliable results.

While we used a specific dataset (EV Charging) to illustrate these steps, the framework itself is general enough to apply to pretty much any data science problem -- whether its predicting customer churn, forecasting sales, detecting fraud, or anything else. The steps remain the same, only the specifics change.

\newpage

% ============================================================
\section*{References}
% ============================================================

\begin{enumerate}
    \item Provost, F., \& Fawcett, T. (2013). \textit{Data Science for Business}. O'Reilly Media.
    \item James, G., Witten, D., Hastie, T., \& Tibshirani, R. (2013). \textit{An Introduction to Statistical Learning}. Springer.
    \item Géron, A. (2019). \textit{Hands-On Machine Learning with Scikit-Learn, Keras, and TensorFlow}. O'Reilly Media.
    \item scikit-learn Documentation. \url{https://scikit-learn.org/stable/}
    \item EV Charging Dataset -- Course provided dataset, FDS Project, VIT University, 2024.
\end{enumerate}

\end{document}
